%---------------------------------------------------------------------
% Part three: cross-surah themes. Three questions.
%---------------------------------------------------------------------

\section*{حصہ سوم --- اہم مضامین}
\markboth{اہم مضامین}{}

%---------------------------------------------------------------------
\begin{sawal}{سوال \textenglish{16} \ \textnormal{\small(خزاں \textenglish{2013}، سوال \textenglish{6})}}
بنی اسرائیل کون تھے؟ ان کی نافرمانیوں کے بارے میں سورۃ البقرہ اور آل عمران
میں بیان شدہ تفصیل کا خلاصہ تحریر کریں۔
\end{sawal}

\begin{hal}
\textbf{تعارف۔} بنی اسرائیل سے مراد حضرت یعقوب\,علیہ السلام (لقب: اسرائیل)
کی اولاد، یعنی موسیٰ و ہارون\,علیہما السلام کی امت۔ قرآن نے انہیں سب سے
زیادہ نصیحت کی --- صرف البقرہ میں تیس سے زائد آیات انہی کے احوال پر ہیں۔

\medskip
\ayat{يَا بَنِي إِسْرَائِيلَ اذْكُرُوا نِعْمَتِيَ الَّتِي أَنْعَمْتُ
عَلَيْكُمْ وَأَوْفُوا بِعَهْدِي أُوفِ بِعَهْدِكُمْ}{البقرہ، \textenglish{40}}
\tarjuma{اے بنی اسرائیل! میری وہ نعمت یاد کرو جو میں نے تم پر کی، اور میرا
عہد پورا کرو، میں تمہارا عہد پورا کروں گا۔}

\medskip
\textbf{نافرمانیاں --- خلاصہ:}
\begin{itemize}\setlength{\itemsep}{1pt}
  \item \textbf{بچھڑے کی پوجا۔} موسیٰ\,علیہ السلام کے طور پر جانے کے تھوڑے
        عرصے بعد ہی گؤسالہ پرستی شروع کر دی (\textenglish{51--54})۔
  \item \textbf{انبیاء کا انکار اور قتل۔} \ar{أَفَكُلَّمَا جَاءَكُمْ رَسُولٌ
        بِمَا لَا تَهْوَىٰ أَنفُسُكُمُ اسْتَكْبَرْتُمْ فَفَرِيقًا كَذَّبْتُمْ
        وَفَرِيقًا تَقْتُلُونَ} --- جب بھی کوئی رسول تمہاری خواہش کے خلاف
        لایا، تکبر کیا، کسی کو جھٹلایا، کسی کو قتل کیا (\textenglish{87})۔
  \item \textbf{عہد شکنی۔} بار بار عہد کر کے توڑا (\textenglish{83--85})۔
  \item \textbf{کتاب میں تحریف اور کتمانِ حق۔} (تفصیل سوال \textenglish{3}
        اور \textenglish{14} میں)۔
  \item \textbf{حسد۔} رسول کریم \saw کی نبوت کو پہچاننے کے باوجود محض حسد سے
        انکار کیا (سوال \textenglish{3})۔
\end{itemize}

\medskip
\textbf{سبق۔} انعام کے ساتھ ذمہ داری بھی آتی ہے --- علم رکھنے کے باوجود
نافرمانی کرنا لاعلمی کی نافرمانی سے زیادہ سنگین ہے۔ یہی وجہ ہے کہ قرآن بنی
اسرائیل کے قصے بار بار دہراتا ہے: تاکہ امتِ مسلمہ انہی غلطیوں سے بچے۔
\end{hal}

%---------------------------------------------------------------------
\begin{sawal}{سوال \textenglish{17} \ \textnormal{\small(بہار \textenglish{2022}، سوال \textenglish{8} اور خزاں \textenglish{2013}، سوال \textenglish{5}، سوال \textenglish{7})}}
سود کی حرمت پر نوٹ لکھیے، نیز انفاق فی سبیل اللہ کی اہمیت آیات کی روشنی میں
بیان کریں۔
\end{sawal}

\begin{hal}
\textbf{(الف) سود کی حرمت۔}

\ayat{يَا أَيُّهَا الَّذِينَ آمَنُوا لَا تَأْكُلُوا الرِّبَا أَضْعَافًا
مُّضَاعَفَةً وَاتَّقُوا اللَّهَ لَعَلَّكُمْ تُفْلِحُونَ}{آل عمران،
\textenglish{130}}
\tarjuma{اے ایمان والو! سود کئی گنا بڑھا چڑھا کر مت کھاؤ، اور اللہ سے ڈرو
تاکہ فلاح پاؤ۔}

\smallskip
\textbf{تشریح:} یہ زمانۂ جاہلیت کے سود کی مذمت ہے، جس میں قرض کی مدت پوری
ہونے پر رقم کئی گنا بڑھا دی جاتی تھی۔ سورۃ البقرہ (\textenglish{275--279})
میں حکم مزید سخت آیا: \ar{وَأَحَلَّ اللَّهُ الْبَيْعَ وَحَرَّمَ الرِّبَا}
--- اللہ نے تجارت حلال اور سود حرام کیا؛ اور \ar{فَإِن لَّمْ تَفْعَلُوا
فَأْذَنُوا بِحَرْبٍ مِّنَ اللَّهِ وَرَسُولِهِ} --- اگر سود نہ چھوڑا تو اللہ
اور اس کے رسول سے جنگ کا اعلان سمجھو (\textenglish{279})۔ \textbf{وجہ حرمت:}
سود دولت کو چند ہاتھوں میں مرتکز کرتا ہے، محنت کے بغیر منافع دیتا ہے، اور
غریب کو مزید غریب کرتا ہے۔

\medskip
\textbf{(ب) انفاق فی سبیل اللہ۔}

\ayat{مَّثَلُ الَّذِينَ يُنفِقُونَ أَمْوَالَهُمْ فِي سَبِيلِ اللَّهِ كَمَثَلِ
حَبَّةٍ أَنبَتَتْ سَبْعَ سَنَابِلَ فِي كُلِّ سُنبُلَةٍ مِّائَةُ حَبَّةٍ}
{البقرہ، \textenglish{261}}
\tarjuma{جو لوگ اپنا مال اللہ کی راہ میں خرچ کرتے ہیں، ان کی مثال ایک دانے کی
سی ہے جس سے سات بالیں نکلیں، ہر بالی میں سو دانے۔}

\smallskip
\textbf{تشریح:} انفاق کا اجر \textbf{سات سو گنا تک} بڑھایا جاتا ہے، اور آگے
فرمایا: \ar{وَاللَّهُ يُضَاعِفُ لِمَن يَشَاءُ} --- اللہ جس کے لیے چاہے مزید
بڑھا دیتا ہے۔ شرط: احسان جتانا اور ایذا دینا اجر ضائع کر دیتا ہے
(\textenglish{262,264})۔ سود اور انفاق کا موازنہ خاص ہے --- ایک مال کو
سمیٹتا ہے، دوسرا بڑھاتا ہے؛ ایک اللہ کی ناراضی، دوسرا اس کی رضا۔
\end{hal}

%---------------------------------------------------------------------
\begin{sawal}{سوال \textenglish{18} \ \textnormal{\small(بہار \textenglish{2013}، سوال \textenglish{3} اور خزاں \textenglish{2013}، سوال \textenglish{8})}}
درج ذیل الفاظ کے معنی لکھیے اور بتائیں کہ یہ کس آیت اور کس سیاق سے تعلق
رکھتے ہیں۔
\end{sawal}

\begin{hal}
دونوں پرچوں میں یہ سوال \textbf{لازمی نوعیت کا} تھا: دس سے پندرہ الفاظ، ہر ایک
کا مختصر جواب۔ سورۃ البقرہ کے چند اہم الفاظ، جو دونوں پرچوں میں پوچھے گئے:

\medskip
\begin{center}
\begin{tabular}{@{}p{2.6cm} p{5.4cm} p{5.6cm}@{}}
\toprule
\textbf{لفظ} & \textbf{معنی} & \textbf{سیاق} \\
\midrule
اَشْہُرٌ مَعْلُومٰتٌ & معلوم و مقررہ مہینے & حج کے تین مہینے --- شوال، ذی
  قعدہ، ذی الحجہ (\textenglish{197}) \\
رَفَث & شہوانی گفتگو یا فعل & حالتِ احرام میں ممنوع (\textenglish{197}) \\
فُسُوق & نافرمانی، گناہ کی بات & حجاج پر ممنوع (\textenglish{197}) \\
جِدَال & جھگڑا، بحث و تکرار & حج کے دوران ممنوع (\textenglish{197}) \\
خُلْع & عورت کا مہر یا مال دے کر
  شوہر سے طلاق لینا & نکاح ختم کرنے کی ایک صورت (\textenglish{229}) \\
لَغْوٌ فِیْ الْاَیْمَان & بلا ارادہ، عادتاً نکل جانے والی قسم & معاف، مؤاخذہ
  نہیں (\textenglish{225}) \\
عُرْضَۃً & نشانہ، بہانہ & اللہ کے نام کو نیکی چھوڑنے کا بہانہ نہ بنانا
  (\textenglish{224}) \\
مُصْلِحُوْن & اصلاح کرنے والے (منافقین کا جھوٹا دعویٰ) & فساد کو اصلاح کہنا
  (\textenglish{11}) \\
سُفَہَاء & بے وقوف، کم عقل & منافقین کا مومنوں کو یہ کہنا (\textenglish{13}) \\
عُرْوَۃُ الْوُثْقٰی & مضبوط کڑا، ٹوٹنے والا نہیں & ایمان باللہ کی تشبیہ
  (\textenglish{256}) \\
\bottomrule
\end{tabular}
\end{center}

\smallskip
\textbf{صیغوں کی پہچان کا اصول:} فعل کا صیغہ پہچاننے کے لیے تین چیزیں دیکھیے
--- \textbf{زمانہ} (ماضی، مضارع، امر)، \textbf{فاعل کی تعداد و صنف} (واحد،
تثنیہ، جمع؛ مذکر، مؤنث)، اور \textbf{باب}۔ مثلاً \ar{يَسْتَبْدِلُونَ} مضارع
جمع مذکر غائب (باب استفعال) ہے۔
\end{hal}